\documentclass[12pt]{article}
\usepackage[T1]{fontenc}
\usepackage[utf8]{inputenc}
\usepackage{lmodern}
\usepackage{textcomp}
\usepackage{enumitem}
\usepackage{geometry}
\geometry{margin=1in}
\begin{document}
\begin{center}
Answer Key - Sociology - Undergraduate Level Assessment 20 \
\small (Answers are ordered exactly as in the question paper.)
\end{center}

\section*{Multiple Choice}
\begin{enumerate}[label=\arabic*.]
\item \textbf{Answer:} A
\\ \textit{Options:} A) we live in a world of superficial, fragmented images; B) no theory is better than any other: 'anything goes'; C) society has changed and we need new kinds of theory; D) all of the above
\\ \textit{Note:} Correct option: A - we live in a world of superficial, fragmented images
\item \textbf{Answer:} D
\\ \textit{Options:} A) the ruling class, or bourgeoisie, who exploit the proletariat; B) a capitalist class dependent on property ownership and advantaged life chances; C) an alignment of classes with shared interests but no state power; D) a state elite whose members are drawn overwhelmingly from a power bloc
\\ \textit{Note:} Correct option: D - a state elite whose members are drawn overwhelmingly from a power bloc
\item \textbf{Answer:} D
\\ \textit{Options:} A) a complex network of interaction at a micro-level; B) a source of conflict, inequality, and alienation; C) an unstable structure of social relations; D) a normative framework of roles and institutions
\\ \textit{Note:} Correct option: D - a normative framework of roles and institutions
\item \textbf{Answer:} D
\\ \textit{Options:} A) Prime Minister, Archbishop of York, Viscounts of England; B) Marquesses of England, Earls of Great Britain, King's Brothers; C) Esquires, Serjeants of Law, Dukes' Eldest Sons; D) King's Grandsons, Lord High Treasurer, Companions of the Bath
\\ \textit{Note:} Correct option: D - King's Grandsons, Lord High Treasurer, Companions of the Bath
\item \textbf{Answer:} B
\\ \textit{Options:} A) the proportion of people living alone has fallen to 29\%; B) many people are cohabiting in long term relationships; C) the upward curve of remarriages compensates for the drop in first marriages; D) all of the above
\\ \textit{Note:} Correct option: B - many people are cohabiting in long term relationships
\end{enumerate}

\section*{True / False}
\begin{enumerate}[label=\arabic*.]
\setcounter{enumi}{5}
\item \textbf{Answer:} False
\\ \textit{Options:} A) True; B) False
\\ \textit{Note:} LLM-generated false statement. Correct answer: 'generous financial benefits for single parents, students and the unemployed'.
\item \textbf{Answer:} False
\\ \textit{Options:} A) True; B) False
\\ \textit{Note:} LLM-generated false statement. Correct answer: 'its knowledge is accumulated from many different research contexts'.
\item \textbf{Answer:} False
\\ \textit{Options:} A) True; B) False
\\ \textit{Note:} LLM-generated false statement. Correct answer: 'the precise, scientific study of observable phenomena'.
\item \textbf{Answer:} True
\\ \textit{Options:} A) True; B) False
\\ \textit{Note:} LLM-generated true statement based on MCQ.
\item \textbf{Answer:} False
\\ \textit{Options:} A) True; B) False
\\ \textit{Note:} LLM-generated false statement. Correct answer: 'individuals buying welfare privately, with some means-tested benefits'.
\end{enumerate}

\section*{Long Form Response}
\begin{enumerate}[label=\arabic*.]
\setcounter{enumi}{10}
\item \textbf{Answer:} supporting LEAs that appeared to be failing. Explain why this is the correct answer and provide supporting evidence. Reference key facts or concepts that support this choice.
\\ \textit{Note:} Long-form derived from MCQ; award credit for stating the correct idea and giving supporting reasoning.
\item \textbf{Answer:} making all strike action illegal. Explain why this is the correct answer and provide supporting evidence. Reference key facts or concepts that support this choice.
\\ \textit{Note:} Long-form derived from MCQ; award credit for stating the correct idea and giving supporting reasoning.
\end{enumerate}

\vfill
\noindent\textit{End of Answer Key}
\end{document}